\section{Related Work}
\label{sec:related}

A survey~\cite{harlow2024survey} reviews mmWave radar in robotic
motion estimation, localization, and mapping.  We focus on
radar-inertial odometry for aerial platforms.

\textbf{Radar ego-velocity estimation.}
Doppler-based ego-velocity estimation from mmWave radar was
demonstrated by Kellner et al.~\cite{kellner2013instantaneous} and
extended to 3D with weighted least squares by
Kramer et al.~\cite{kramer2020rve}, whose RVE system fuses radar
with \ac{imu} in a discrete-time \ac{ekf}.  Doer and
Trommer~\cite{doer2020ekf} present an \ac{ekf}-based \ac{rio} with online
extrinsic calibration.  Our work uses per-frame \ac{wls}
ego-velocity estimates to initialize the position trajectory.
The continuous-time sliding-window optimizer uses one Doppler factor
per return (\autoref{sec:init}, \autoref{sec:robust}).

\textbf{Continuous-time trajectory estimation.}
Furgale et al.~\cite{furgale2013unified} introduced uniform B-splines
for continuous-time visual-inertial calibration.  The basalt
\ac{vio}~\cite{usenko2020basalt} uses cumulative B-splines on Lie groups
for stereo-inertial odometry; we adapt this basalt-style
cumulative SO(3) B-spline to radar-inertial fusion.

\textbf{Marginalization in sliding-window estimators.}
OKVIS~\cite{leutenegger2015keyframe} and
VINS-Mono~\cite{qin2018vins} use Schur complement
marginalization for visual-inertial systems, propagating information
from dropped keyframes as a dense Gaussian prior. We apply the same
principle to continuous-time B-spline control points.

\textbf{Continuous-time radar-inertial odometry.}
Ng et al.~\cite{ng2021ctrio} formulate continuous-time radar-inertial
odometry by fusing Doppler measurements from
\emph{multiple automotive} radars with an \ac{imu} as a least-squares
problem over a B-spline trajectory, evaluated on passenger vehicles.
Burnett et al.~\cite{burnett2025gp} use a Gaussian-process motion
prior for spinning-radar- and lidar-inertial odometry, and
\v{S}tironja et al.~\cite{stironja2026rio} model the \ac{imu} trajectory
continuously inside an \ac{rio} estimator with online spatio-temporal
calibration.  Yang et al.~\cite{yang2025gorio} (Go-RIO)
similarly fuse radar velocity with an \ac{imu} in continuous time and
target the same radar elevation drift with a ground-plane model, which
requires a persistently visible floor.  Our aerobatic flights see it
only intermittently.  Our front end rejects Doppler outliers but does
not impose a height constraint.  Huang et al.~\cite{huang2024lessismore}
exploit radar cross-section and Doppler physics to make a sparse
single-chip cloud usable in a discrete-time sliding-window
optimization, evaluated on ground-robot and handheld platforms.
We study a single-chip radar (10--60 points per frame,
12-channel array) on a quadrotor in aggressive racing and
\rev{aerobatic backflip} flight up to $10$\,rad/s.
The radar measurement errors increase in this regime
(\autoref{sec:backflips}).

\textbf{Single-chip \ac{rio} on aerial platforms.}
Single-chip radar-inertial fusion has also been evaluated on multirotors.
Michalczyk et al.~\cite{michalczyk2022ekf} use a
discrete-time \ac{ekf} with scan-to-scan point matching and Doppler (a
factor-graph reformulation follows in~\cite{michalczyk2024fg}); Girod et
al.~\cite{girod2024brio} run a real-time onboard M-estimator smoother on a
quadrotor in cities and forests, but cap velocity at
$2$\,m/s and anchor the vertical channel with a
barometer.  Our evaluation focuses on continuous-time estimation during
aggressive and \rev{aerobatic} flight, without a barometric height
measurement (\autoref{sec:baselines}).  None of the systems above
covers that combination.

\textbf{Sliding-window marginalization and consistency.}
Sliding-window marginalization can introduce inconsistency
in discrete time~\cite{dongsi2012consistency}, where \ac{fej}
methods fix linearization points to preserve the observability nullspace.
In continuous time, CLIC~\cite{lv2023clic} and
LIO-MARS~\cite{quenzel2025liomars} marginalize spline control points leaving the
window by adjacency to the dropped control points.  We follow the same
construction and account for the globally shared
\ac{imu} bias, which couples every \ac{imu} factor of the window and
requires care when selecting factors for marginalization
(\autoref{sec:sw}).

\textbf{Positioning against the \ac{ekf}-\ac{rio} baselines.}
In the discrete-time UAV line, the \ac{ekf}-\ac{rio} family of Doer and
Trommer~\cite{doer2020ekf} provides our comparison baselines.  Absolute
\ac{rmse} is not comparable across papers (firmware configuration alone
changes the Doppler resolution by an order of magnitude), so we
compare the systems on shared data:
\autoref{sec:baselines} runs our system on their public ICINS-2021
flights~\cite{doer2021yrio} and re-evaluates their estimators under
our causal metrics on the same data.
